% The half-chain collapse for the coupled slice (paper 13 of the O lane).
% Lane: O (operator bridge), module YangMills/OS/SpatialOS.lean.
% REGIME: tricotomy; NO scores, NO desk language, NO commit chronology beyond
% the reproducibility freeze.
\documentclass[11pt]{article}
\usepackage[margin=1.1in]{geometry}
\usepackage{amsmath,amssymb,amsthm}
\usepackage{booktabs,longtable,array}
\usepackage[colorlinks=true,linkcolor=blue,urlcolor=blue,citecolor=blue]{hyperref}

\newtheorem{theorem}{Theorem}[section]
\newtheorem{lemma}[theorem]{Lemma}
\newtheorem{proposition}[theorem]{Proposition}
\newtheorem{corollary}[theorem]{Corollary}
\newtheorem{remark}[theorem]{Remark}
\theoremstyle{definition}
\newtheorem{definition}[theorem]{Definition}

\emergencystretch=3em
\tolerance=2000

\newcommand{\anchor}{56b3d5ef315f13670494fb9e13669ae2d1975554}
\newcommand{\osline}[3]{%
  \href{https://github.com/lluiseriksson/THE-ERIKSSON-PROGRAMME/blob/\anchor/YangMills/OS/#1.lean\#L#2}{\nolinkurl{#3}}}
\newcommand{\repofile}[2]{%
  \href{https://github.com/lluiseriksson/THE-ERIKSSON-PROGRAMME/blob/\anchor/#1}{\nolinkurl{#2}}}
\newcommand{\R}{\mathbb{R}}
\newcommand{\C}{\mathbb{C}}
\newcommand{\Z}{\mathbb{Z}}
\newcommand{\SU}{\mathrm{SU}}

\title{The Collapse:\\
A Machine-Checked Half-Chain Reflection Form\\
for the Coupled $\Z_2$ Slice\\
\large complex observables of a whole half --- two reflections,
and the one link that is verified but not yet proved}
\author{Lluis Eriksson}
\date{}

\begin{document}
\maketitle

\begin{abstract}
\textbf{What is proved, stated before anything else.}  For a complex-valued
observable $F$ of an \emph{entire past half-path} of the coupled $\Z_2$ slice,
the reflected form $\langle \Theta F, F\rangle$ is a non-negative real: through
a \emph{site} for every coupling $\beta$, and through a \emph{bond} for
$\beta\ge0$.  That hypothesis is proved \emph{sufficient} and is not proved
necessary here --- and at $L=0$ it could not be, since the bare kernel is then
the scalar $1$ and the bond form is non-negative at every $\beta$.  The witness
that would make the boundary exact from one site upwards is authorised by a gate
and is not written (\S\ref{sec:open}).  The mechanism is a single map --- summing out the interior of a
half sends $F$ to a vector indexed by its boundary slice alone --- after which
both cases reduce to kernel facts already machine-checked in the companion paper
\cite{ErikssonRefl}.

\textbf{And what is NOT proved, stated second.}  The form is \emph{defined} here
as a sum over pairs of halves.  The lemma identifying it with the Gibbs sum over
whole paths --- that the two halves and the crossing factor assemble into one
path, bijectively --- is \emph{verified} to $10^{-12}$ against a brute-force
enumeration of every path, and is \emph{not proved}.  Until it is, this paper
does not establish the Osterwalder--Schrader axiom; it establishes positivity of
a form that the numerics say is that pairing.  The distinction is the whole
reason this sentence is in the abstract and not in a footnote.

\textbf{Where the two reflections differ, and it is not cosmetic.}  Through a
bond the two halves are disjoint and meet through one kernel factor, so
positivity needs the kernel itself positive semidefinite --- which holds exactly
for $\beta\ge0$.  Through a site the halves \emph{share} the middle slice; the
form is then a sum of squared moduli divided by the weight of that slice, so it
is non-negative for \emph{every} $\beta$, negative coupling included, and no
property of the kernel is used at all.

\textbf{What the source weight does.}  In the companion paper the weight was
handled by congruence: conjugation by $\sqrt{w}$ cannot change the sign of a
quadratic form.  Here it is not conjugated away, it is \emph{summed} away: the
weight of every interior slice is absorbed into the collapse, and what survives
at the boundary is the kernel alone.  Two different reasons, same conclusion ---
nothing in the hypotheses depends on $w$ beyond positivity.

\textbf{What this is not.}  No reconstruction: the physical Hilbert space as the
quotient of the past algebra by the null space of this form is not built.
Nothing here concerns uniformity in the extent, $\SU(N)$, the continuum limit,
or the Yang--Mills mass gap \cite{Clay}.

\textbf{On the pre-registration.}  Four gates were committed before a line of
the module was written; \S\ref{sec:gates} reports what happened to each.  Two of
them read a minimum eigenvalue rather than sampling observables, which is the
instrument the previous campaign's autopsy said was needed.
\end{abstract}

\section{The gap this paper attacks}\label{sec:gap}

The companion paper proves the reflected two-point form of the Gibbs measure
non-negative --- but for \emph{real} observables of a \emph{single} slice,
evaluated at the two \emph{ends} of a path \cite{ErikssonRefl}.  The
Osterwalder--Schrader axiom asks for something wider in three directions at
once: the observable should be a function of the whole past half-chain, it
should be complex, and the statement should be about a matrix, not a number
\cite{OS73,OS75}.  Reflection positivity is one axiom among several, and on a
lattice it is the one that yields a self-adjoint transfer operator by the
standard reconstruction \cite{OSeiler,GJ}.

That paper also said, in its scope section, that the missing step is a
\emph{construction} and not a further inequality: collapse each half to a
boundary vector and the two inequalities you need are the ones already proved.
This paper carries out that construction and reports how far it gets.

\paragraph{All three directions are closed here; the identification is not.}
The observables are now complex, they depend on the entire half-path, and the
statement is about a Gram matrix (Theorem~\ref{thm:gram}).  What remains open is
not one of the three directions but the identification of the form with the path
measure (\S\ref{sec:open}), and that is the one that decides what the theorems
are about.

\section{The objects}\label{sec:objects}

A slice is $\sigma:\{1,\dots,L\}\to\Z_2$ and the decoupled kernel is
\[
  K(\sigma,\tau)\;=\;\prod_{j} e^{\beta\, s(\sigma_j,\tau_j)},
  \qquad s=+1 \text{ if equal},\ -1 \text{ otherwise}.
\]
A \emph{past half-path} is a map $a:\{0,\dots,m\}\to\{\text{slices}\}$, and an
observable of the past half-chain is a function $F$ of one, valued in $\C$.  Its
weight is the ordinary Gibbs weight of that half,
\[
  W(a)\;=\;\Bigl(\prod_{t=0}^{m} w(a_t)\Bigr)
           \Bigl(\prod_{t=0}^{m-1} K(a_t,a_{t+1})\Bigr),
\]
which is the same function the bridge module already uses for whole paths
\cite{ErikssonGibbs}; \textbf{no new notion of weight is introduced anywhere in
this work}.  The reflection $\Theta$ is path reversal.

\begin{definition}[the collapse]\label{def:collapse}
$\displaystyle (\Phi F)(\sigma)\;=\;\sum_{a\,:\,a_m=\sigma} W(a)\,F(a).$
\end{definition}

\osline{SpatialOS}{82}{collapse}, with its linearity at
\osline{SpatialOS}{87}{collapse\_add} and
\osline{SpatialOS}{96}{collapse\_smul}.  The regrouping of a sum over halves by
boundary slice --- the only combinatorial fact either factorisation needs --- is
\osline{SpatialOS}{105}{sum\_halves\_by\_edge}.

\section{Through a site: every $\beta$}\label{sec:site}

With the plane through slice $m$ the two halves share that slice.  Its weight is
carried by both, so it is counted once.

\begin{theorem}[site factorisation]\label{thm:site}
$\displaystyle \langle\Theta F,F\rangle_{\mathrm{site}}
   \;=\;\sum_\sigma \frac{|(\Phi F)(\sigma)|^{2}}{w(\sigma)}.$
\end{theorem}

\osline{SpatialOS}{142}{osPairingSite\_eq}.  Nothing about $K$ enters: the
halves communicate only through the shared slice.

\begin{corollary}\label{cor:site}
For every strictly positive $w$ and \emph{every} $\beta$ --- negative coupling
included --- the form is a non-negative real.
\end{corollary}

\osline{SpatialOS}{170}{osPairingSite\_nonneg}.

\section{Real to complex, once and generically}\label{sec:complex}

\begin{lemma}\label{lem:complex}
For a real \emph{symmetric} kernel,
$\sum_{i,j}\overline{z_i}K_{ij}z_j
  = \sum_{i,j}K_{ij}\,x_ix_j + \sum_{i,j}K_{ij}\,y_iy_j$
with $z=x+iy$.
\end{lemma}

\osline{SpatialOS}{200}{complexQuad\_eq}, resting on the cancellation of the
cross terms at
\osline{SpatialOS}{189}{sum\_cross\_antisymm}.  This is the entire content of
``real observables to complex ones'', and it costs exactly symmetry.  It is
stated for an arbitrary finite index type, so it is reusable outside this lane.

\section{Through a bond: sufficiency of $\beta\ge0$}\label{sec:bond}

\begin{theorem}[bond factorisation]\label{thm:bond}
$\displaystyle \langle\Theta F,F\rangle_{\mathrm{bond}}
   \;=\;\sum_{\sigma,\tau}\overline{(\Phi F)(\sigma)}\,K(\sigma,\tau)\,
       (\Phi F)(\tau).$
\end{theorem}

\osline{SpatialOS}{244}{osPairingBond\_eq}.  The source weight is not conjugated
away; it is summed away, into $\Phi$.  What is left at the boundary is the bare
kernel.

\begin{theorem}[the form, for $\beta\ge0$]\label{thm:bondpos}
For every strictly positive $w$ and every $\beta\ge0$, $\langle\Theta
F,F\rangle_{\mathrm{bond}}$ is a non-negative real, for every complex observable
$F$ of the whole past half-chain.
\end{theorem}

\osline{SpatialOS}{297}{osPairingBond\_nonneg}, which is
Theorem~\ref{thm:bond} composed with Lemma~\ref{lem:complex} and the companion
paper's positive semidefiniteness of $K$ at $\beta\ge0$ \cite{ErikssonRefl}.

\paragraph{The matrix, not one diagonal entry.}  The axiom is about a Gram
matrix.  The cross form $\langle\Theta F,G\rangle_{\mathrm{bond}}$ is defined
separately, factors through the collapse by the same argument, and the matrix is
assembled from it.

\begin{theorem}[the Gram identity]\label{thm:gram}
For every finite family $F_1,\dots,F_k$ of complex half-chain observables and
every complex $c_1,\dots,c_k$,
\[
  \sum_{i,j}\overline{c_i}\,c_j\,\langle\Theta F_i,F_j\rangle_{\mathrm{bond}}
  \;=\;\Bigl\langle \Theta\Bigl(\sum_i c_i F_i\Bigr),
                    \sum_i c_i F_i\Bigr\rangle_{\mathrm{bond}},
\]
and for $\beta\ge0$ both sides are the same non-negative real.
\end{theorem}

\osline{SpatialOS}{427}{osPairingBond\_gram} and
\osline{SpatialOS}{456}{osPairingBond\_gram\_nonneg}, with the cross form at
\osline{SpatialOS}{349}{osPairingBondCross}, its factorisation at
\osline{SpatialOS}{375}{osPairingBondCross\_eq}, the linearity of the collapse at
\osline{SpatialOS}{362}{collapse\_sum}, and the four-fold reindexing the assembly
needs stated on its own at \osline{SpatialOS}{330}{sum\_comm4}.  Positivity is
\emph{transported} to the matrix rather than re-proved: the identity says the
Gram sum \emph{is} the form of the combined observable.

\section{The four gates}\label{sec:gates}

\repofile{scripts/judge\_spatial\_os.py}{scripts/judge\_spatial\_os.py},
committed before a line of the module existed.

The plan for this campaign announced \emph{three} gates.  Writing them down
showed that the site collapse and the bond collapse are two different identities
licensing two different theorems, and this lane has twice been punished for
putting two claims behind one verdict.  The count was moved to four, in public,
before anything was run.  Raising a pre-registration before it executes is
legitimate; discovering afterwards that one gate was really two is not.

\begin{itemize}
\item \textbf{A1, A2} --- the two collapse identities, as \emph{matrix}
  identities to $10^{-12}$, against a pairing matrix built by brute-force
  enumeration of every path, so that they do not presuppose the factorisation
  they test.  Both \textsc{pass}.
\item \textbf{B} --- the \emph{minimum eigenvalue} of that matrix at
  $\beta\ge0$.  The previous campaign's autopsy measured that drawing random
  observables can miss a violating direction carrying $10^{-3}$ of the weight;
  an eigenvalue cannot miss it.  \textsc{pass}.
\item \textbf{C} --- below zero at odd separation the minimum eigenvalue must be
  strictly negative, and in the cell $(L=1,m=0,w\equiv1)$ it must equal
  $e^{\beta}-e^{-\beta}$ to $10^{-12}$.  Measured agreement: $\le2.2\times
  10^{-16}$.  \textsc{pass}.
\end{itemize}

Gate C licenses a sharpness witness that this paper does not yet contain
(\S\ref{sec:open}).  A passed gate is a permission, not a theorem.

\section{What is open, and which part is load-bearing}\label{sec:open}

\paragraph{The link that decides what this paper is.}  The two forms are
\emph{defined} as sums over pairs of halves.  That this equals the Gibbs sum
over whole paths requires the assembly map --- past half, crossing factor,
reversed future half --- to be a bijection onto paths, and the half weights to
multiply to the path weight.  Gates A1 and A2 verify exactly this, to
$10^{-12}$, against brute-force enumeration; it is not proved.  \textbf{Until it
is, the theorems above are positivity statements about a form the numerics
identify with the Osterwalder--Schrader pairing, and not about that pairing.}
This is the first thing a reader should check when the next version appears.

\paragraph{The sharpness witness.}  Gate C authorises it and it is not written.
Consequently $\beta\ge0$ appears above as a sufficient hypothesis only.  The
route is identified: at $m=0$ the collapse is multiplication by $w$, hence
surjective when $w>0$, so an observable can be chosen whose boundary vector is
the sign observable that the companion paper already drives negative below zero
\cite{ErikssonRefl}.

\paragraph{No reconstruction.}  The quotient by the null space of this form,
with the transfer operator descending to a self-adjoint contraction, is not
built.  Nothing here concerns uniformity in the extent, $\SU(N)$, the continuum
limit, or the Clay problem \cite{Clay}.

\section{Reproducibility}\label{sec:repro}

All artifacts are at source checkpoint \texttt{\anchor} on branch \texttt{main}.
Every hyperlink was resolved against that commit, and checked to point at an
actual declaration line, before compilation.

\begin{center}
\begin{tabular}{ll}
\toprule
Quantity & Value \\
\midrule
Full core build & 8464 jobs, success \\
\quad checkpoint before this module & 8463 jobs \\
Oracle commands, and answers & 2773 \\
\quad distinct declarations & 2770 \\
\quad commands listed twice & 3 \\
\quad distinct, with axiom dependencies & 2746 \\
\quad distinct, axiom-free & 24 \\
Declarations in this module & 21 \\
\quad of them in the oracle & 21 \\
Distinct axioms across all outputs & \texttt{propext}, \texttt{Quot.sound},
  \texttt{Classical.choice} \\
Occurrences of \texttt{sorryAx} & 0 \\
Project axioms & 0 \\
Errors & 0 \\
\bottomrule
\end{tabular}
\end{center}

The job count incremented from 8463 to 8464 when the module joined the core,
which is this repository's standing check that a new module is really being
elaborated.  It did not move again when declarations were added to an existing
module, which is the correct reading of that check rather than a failure of it.
The oracle counters close: $2746+24=2770$ and $2770+3=2773$, the three repeats
being long-standing duplicate lines elsewhere in the file.  Every declaration of
this module appears in the oracle, reconciled by name.

The module is \repofile{YangMills/OS/SpatialOS.lean}{YangMills/OS/SpatialOS.lean}
and the companion is
\repofile{YangMills/OS/SpatialReflection.lean}{YangMills/OS/SpatialReflection.lean}.

\begin{thebibliography}{9}
\bibitem{OS73} K. Osterwalder and R. Schrader,
  \emph{Axioms for Euclidean Green's functions},
  Comm. Math. Phys. \textbf{31} (1973), 83--112.
\bibitem{OS75} K. Osterwalder and R. Schrader,
  \emph{Axioms for Euclidean Green's functions II},
  Comm. Math. Phys. \textbf{42} (1975), 281--305.
\bibitem{OSeiler} K. Osterwalder and E. Seiler,
  \emph{Gauge field theories on a lattice},
  Ann. Physics \textbf{110} (1978), 440--471.
\bibitem{GJ} J. Glimm and A. Jaffe,
  \emph{Quantum Physics: A Functional Integral Point of View},
  2nd ed., Springer, 1987.
\bibitem{ErikssonRefl} L. Eriksson,
  \emph{The Weight That Could Not Break It: Machine-Checked Endpoint Reflection
  Positivity for the Coupled $\Z_2$ Slice}, 2026.
\bibitem{ErikssonGibbs} L. Eriksson,
  \emph{The Measure the Spectral Results Were About}, 2026.
\bibitem{Clay} A. Jaffe and E. Witten,
  \emph{Quantum Yang--Mills Theory}, Clay Mathematics Institute, 2000.
\end{thebibliography}

\end{document}
